% TOP_PROBLEM: {"id":30007759,"problem_number":"LOCAL-30007759","title":"Efficient tax enforcement","category_id":21,"category":{"id":21,"name":"economics","display_name":"Economics","description":"Open economic theory statements, published research agendas, and proposed scoped specifications; question provenance and historical priority are qualified per record.","slug":"economics","order_index":21,"created_at":"2026-10-08T00:00:00Z"},"status":"provenance_pending","external_url":null,"legacy_ids":[],"published":false,"statement":"In the explicitly specified setting: Which mix of audits, third-party reporting, simplification, and social norms delivers the greatest durable increase in tax compliance? The mathematical statement above fixes the target, assumptions and scope of this catalogue formulation.","economics":{"original_id":248,"field":"Taxation and social insurance","kind":"design","formalization_basis":"proposed_specification","source_status":"not_verified","priority_status":"not_established","priority_note":"An original explicit open-question statement was not verified; the cited supporting paper must not be treated as its first listing.","earliest_verified_reference":null,"current_reference":{"authors":["Joel Slemrod"],"title":"Tax Compliance and Enforcement","year":2019,"url":"https://www.aeaweb.org/articles?id=10.1257/jel.20181437","doi":"10.1257/jel.20181437","locator":"Journal of Economic Literature 57(4), 904–954; publisher abstract","evidence_summary":"Reviews audits, information reporting, remittance and experimental methods. The accessible abstract does not state the catalogue long-run joint-policy optimization as an open problem."},"formalization_note":"Catalogue-authored scoped formalization. Population, instruments, welfare objective and identification restrictions must be instantiated before claiming a resolution; source authors are not credited with these added assumptions."},"statusQualification":"Provenance pending: an explicit published statement of this exact open question was not verified. This is a catalogue-authored candidate specification, not a certified open conjecture. Catalogue-authored scoped formalization. Population, instruments, welfare objective and identification restrictions must be instantiated before claiming a resolution; source authors are not credited with these added assumptions.","statusReviewedAt":"2026-10-08"}
% FIRST_POSTED: 2026-10-08
% ENTRY_KIND: statement_only
% !TeX program = lualatex
\documentclass[11pt]{article}
\usepackage{fontspec}
\usepackage[a4paper,margin=25mm]{geometry}
\usepackage{amsmath,amssymb,mathtools}
\usepackage{xurl}
\usepackage[hidelinks]{hyperref}
\newcommand{\sref}[2]{\hyperlink{source:#1:#2}{[S#2]}}
\setlength{\emergencystretch}{3em}
\begin{document}
% problemId:30007759
\section*{Efficient tax enforcement}\label{30007759}
\subsection{Definitions and mathematical statement}
Fix a tax jurisdiction, eligible taxpayers, horizon \(H\), and policy budget \(B\). An enforcement policy \(a=(a_1,a_2,a_3,a_4)\) specifies audit assignment, third-party reporting coverage, filing simplification and communication about social norms, including intensity and timing. Let \(\mathcal A_B\) be a prespecified admissible set whose discounted administrative cost is at most \(B\). For taxpayer \(i\), let \(L_{it}(a)\) be true statutory liability and \(P_{it}(a)\) correctly remitted tax in period \(t\), with fines recorded separately. Define the collection gap and durable policy value by
\[G_H(a)=\mathbb E\!\left[\sum_{t=0}^H\beta^t\sum_i\{L_{it}(a)-P_{it}(a)\}\right],\quad a^*\in\arg\min_{a\in\mathcal A_B}G_H(a),\]
where \(0<\beta\leq1\) is a specified discount factor. Identify or bound policy interactions and post-intervention adaptation using randomized assignment or a defensible quasi-experiment, allowing interference through taxpayer and adviser networks. Define liability measurement from representative verification audits or another justified observation model. Also report real activity and compliance costs, because smaller measured gaps may arise from changed liability rather than improved honesty. Determine which combinations dominate isolated interventions and whether rankings persist after monitoring ends. The admissible set and adaptation horizon are part of the proposed specification.

\par\textbf{Formulation and scope.} Catalogue-authored scoped formalization. Population, instruments, welfare objective and identification restrictions must be instantiated before claiming a resolution; source authors are not credited with these added assumptions.

\par\textbf{Open-question provenance.} An explicit published open-question statement was not verified. Sources below are supporting literature and must not be treated as a first listing.

\par\textbf{Historical priority.} An original explicit open-question statement was not verified; the cited supporting paper must not be treated as its first listing.

\subsection{Short English statement}
In the explicitly specified setting: Which mix of audits, third-party reporting, simplification, and social norms delivers the greatest durable increase in tax compliance? The mathematical statement above fixes the target, assumptions and scope of this catalogue formulation.

\subsection{Sources}
\begin{itemize}\raggedright
\item[\textnormal{[S1]}]\hypertarget{source:30007759:1}{} Joel Slemrod. 2019. Tax Compliance and Enforcement. Journal of Economic Literature 57(4), 904–954; publisher abstract.
\url{https://www.aeaweb.org/articles?id=10.1257/jel.20181437}.
DOI: 10.1257/jel.20181437.
{\small\textit{Supporting or current literature; see evidence qualification. Reviews audits, information reporting, remittance and experimental methods. The accessible abstract does not state the catalogue long-run joint-policy optimization as an open problem.}}
\end{itemize}
\end{document}
